\hbadness=10000
%==============================================================================
% Homework Assignment Template
%==============================================================================
\documentclass[11pt]{article}

%------------------------------------------------------------------------------
% PACKAGES
%------------------------------------------------------------------------------
\usepackage[margin=1in]{geometry}   % Adjust margins
\usepackage{amsmath,amssymb,amsthm} % Math packages
\usepackage{graphicx}               % Include images
\usepackage{hyperref}               % Hyperlinks
\usepackage{enumitem}               % Better control over lists
\usepackage{fancyhdr}               % Header/footer customization

%------------------------------------------------------------------------------
% CUSTOM COMMANDS & SETTINGS
%------------------------------------------------------------------------------
\newcommand{\courseName}{Ensemble Quantum Computing Laboratory}     % Change to your course name
\newcommand{\assignmentTitle}{Prelab 3} % Change to your assignment title
\newcommand{\studentName}{Runzhao Guo}       % Change to your name
\newcommand{\Date}{Feb 17, 2025}   % Change to your due date

% Set up the page header
\pagestyle{fancy}
\fancyhf{} % Clear header and footer
\lhead{\courseName}
\chead{\assignmentTitle}
\rhead{\studentName}
\rfoot{\thepage}



%------------------------------------------------------------------------------
% DOCUMENT START
%------------------------------------------------------------------------------
\begin{document}
%------------------------------------------------------------------------------
% TITLE SECTION
%------------------------------------------------------------------------------
\begin{center}
    {\large \textbf{\courseName}} \\
    \vspace{3pt}
    {\Large \textbf{\assignmentTitle}} \\
    \vspace{3pt}
    \studentName \\
    \vspace{3pt}
    \Date
\end{center}

\vspace{0.5in} % Add some space before the content starts

\section*{Problem 1}
\subsection*{part a).}

The spectra is shown in Figure \ref{fig:prob1a}.

\begin{figure}[h]
    \centering
    \includegraphics[width=0.5\textwidth]{/Users/runzhaoguo/Documents/MQST-UCLA/411/lab3/figure/prob1_a.jpeg}
    \caption{the spectra for the four cases when only one of a, b, c, and d is nonzero}
    \label{fig:prob1a}
\end{figure}

\subsection*{part b).}
The integral of the given thermal equilibrium is shown in Figure \ref{fig:prob1b}. Notice the frequence increase from left to right.

\begin{figure}[h]
    \centering
    \includegraphics[width=0.5\textwidth]{/Users/runzhaoguo/Documents/MQST-UCLA/411/lab3/figure/prob1_b.png}
    \caption{the integral for given thermal equilibrium density matrix}
    \label{fig:prob1b}
\end{figure}

\subsection*{part c).}

For CNOT gate, the density matrix is processed as such:$\rho=U_{c n} \rho_{\text {therm }} U_{c n}^{\dagger}$, the result we get is: 
\[
\begin{pmatrix}
1/4 + 5\times10^{-4} & 0 & 0 & 0 \\
0 & 1/4 + 3\times10^{-4} & 0 & 0 \\
0 & 0 & 1/4 - 5\times10^{-4} & 0 \\
0 & 0 & 0 & 1/4 - 3\times10^{-4} \\
\end{pmatrix}
\]

For near-CNOT gate, the density matrix is processed as such:$\rho=U_{n c n} \rho_{\text {therm }} U_{n c n}^{\dagger}$, the result we get is:
\[
\begin{pmatrix}
1/4 - i5\times10^{-4} & 0 & 0 & 0 \\
0 & 1/4 + i3\times10^{-4} & 0 & 0 \\
0 & 0 & 1/4 + 5\times10^{-4} & 0 \\
0 & 0 & 0 & 1/4 + 3\times10^{-4} \\
\end{pmatrix}
\]

\begin{figure}[h]
    \centering
    \includegraphics[width=0.5\textwidth]{/Users/runzhaoguo/Documents/MQST-UCLA/411/lab3/figure/prob1_c1.jpeg}
    \caption{the integral for given thermal equilibrium density matrix after CNOT gate}
    \label{fig:prob1c1}
\end{figure}

\begin{figure}[h]
    \centering
    \includegraphics[width=0.5\textwidth]{/Users/runzhaoguo/Documents/MQST-UCLA/411/lab3/figure/prob1_c2.jpeg}
    \caption{the integral for given thermal equilibrium density matrix after near-CNOT gate}
    \label{fig:prob1c1}
\end{figure}

The reletive relation between the integral of two peeks for Proton remain the same.

\section*{Problem 2}

\subsection*{part a).}
For a given Hamiltonian, the thermal equilibrium density matrix $\rho_{eq}$ is given by:
\[
    \rho_{eq} = \frac{e^{-\beta H}}{Tr[e^{-\beta H}]}\\
\]
where $\beta = \frac{1}{K_B T}$

The Hamiltonian for this two-qubit system is $\hat{H} = -\hbar \omega_{H}(1+\delta_H) \hat{I}_Z -\hbar \omega_{C}(1+\delta_C) \hat{S}_Z + 2\pi J\hat{I}_Z \hat{S}_Z$

Thank goodness it's a diagonal matrix, so the exponential can be calculated directly:
\[
    \rho_{eq} = \frac{1}{\sum_{i=1}^{4}e^{-\beta E_i}}
    \begin{pmatrix}
        e^{-\beta E_1} & 0 & 0 & 0 \\
        0 & e^{-\beta E_2} & 0 & 0 \\
        0 & 0 & e^{-\beta E_3} & 0 \\
        0 & 0 & 0 & e^{-\beta E_4} \\
    \end{pmatrix}
\]
Where:\\
    $E_1 = -0.5\hbar (\omega_{H}(1+\delta_H) + \omega_{C}(1+\delta_C) - \pi J)$ \\
    $E_2 = -0.5\hbar (\omega_{H}(1+\delta_H) - \omega_{C}(1+\delta_C) + \pi J)$ \\
    $E_3 = +0.5\hbar (\omega_{H}(1+\delta_H) - \omega_{C}(1+\delta_C) - \pi J)$ \\
    $E_4 = +0.5\hbar (\omega_{H}(1+\delta_H) + \omega_{C}(1+\delta_C) + \pi J)$ \\
After calculation, the thermal equilibrium density matrix is:
\[
    \frac{1}{4.000 + 5.66\times10^{-11}}
    \begin{pmatrix}
    1+6.306\times10^{-6} & 0 & 0 & 0 \\
    0 & 1+3.771\times10^{-6} & 0 & 0 \\
    0 & 0 & 1-3.771\times10^{-6} & 0 \\
    0 & 0 & 0 & 1-6.306\times10^{-6} \\
    \end{pmatrix}
\]

\subsection*{part b}
In my understanding, expressed in product operators, the density matrix is:
\[
    0.25\hat{I_I}\hat{S_I} + 0.5\times10^{-5}\hat{I_I}\hat{S_Z} + 0.13\times10^{-5}\hat{I_Z}\hat{S_I}
\]

Expressing it in a similar way used in Knill paper: 
\[
    \frac{1}{4.000}
    \begin{pmatrix}
    1 & 0 & 0 & 0 \\
    0 & 1 & 0 & 0 \\
    0 & 0 & 1 & 0 \\
    0 & 0 & 0 & 1 \\
    \end{pmatrix}
    + 10^{-5}
    \begin{pmatrix}
        0.157 & 0 & 0 & 0 \\
        0 & 0.094 & 0 & 0 \\
        0 & 0 & -0.094 & 0 \\
        0 & 0 & 0 & -0.157 \\
    \end{pmatrix}
\]

In comparison, the second element in my density matrix is 1/6 of that in the Knill paper. One possible reason is The magnetic field used in the Knill paper is stronger than that in my case, which leads to a larger value of $\omega_{H/C}$.


\section*{Problem 3}
The exaustive average is:
\[
    \begin{pmatrix}
    0.25+1.57\times10^{-6} & 0 & 0 & 0 \\
    0 & 0.25-3.2\times10^{-7} & 0 & 0 \\
    0 & 0 & 0.25-3.2\times10^{-7} & 0 \\
    0 & 0 & 0 & 0.25-3.2\times10^{-7} \\
    \end{pmatrix}
\]
The excess probability is $\rho_{00} - \bar{p} = 1.89\times10^{-6}$

\section*{Problem 4}

\subsection*{part a).}
\[
\mathrm{CNOT}_{(1\to 2)} 
= \frac{1}{2} \Bigl( |0\rangle \langle0| \otimes I + |1\rangle \langle1| \otimes X \Bigr).
\]

\[
    |0\rangle \langle0| = \frac{1}{2} (I + Z)
\]

\[
    |1\rangle \langle1| = \frac{1}{2} (I - Z)
\]

\[
\mathrm{CNOT}_{(1\to 2)} 
= \frac{1}{2} \Bigl( I \otimes I \;+\; Z \otimes I \;+\; I \otimes X \;-\; Z \otimes X \Bigr).
\]
\[
=\frac{1}{4}
    \begin{pmatrix}
    1 + 1  & 1 - 1 & 0 & 0 \\
    1 - 1 & 1 + 1  & 0 & 0 \\
    0 & 0 & 1 - 1 & 1 + 1 \\
    0 & 0 & 1 + 1 & 1 - 1 \\
    \end{pmatrix}
\]

\[
=\frac{1}{2}
    \begin{pmatrix}
    1 & 0 & 0 & 0 \\
    0 & 1 & 0 & 0 \\
    0 & 0 & 0 & 1 \\
    0 & 0 & 1 & 0 \\
    \end{pmatrix}
\]

\subsection*{part b).}


% \begin{thebibliography}{9}
% \bibitem{exampleRef} Author Name, \textit{Title of the Book}, Publisher, Year.
% \end{thebibliography}
%------------------------------------------------------------------------------
% DOCUMENT END
%------------------------------------------------------------------------------
\end{document}
